\subsection{Architecture du pipeline}
Pour rendre le processus clair et compréhensible, nous avons divisé le pipeline en étapes illustrées comme suit : 

\begin{enumerate}
    \item \textbf{Encoder le texte cible} : en utilisant SONAR pour obtenir un vecteur d'embedding de taille 1024 dimensions. Cette étape permet d'extraire des caractéristiques sémantiques et linguistiques qui seront interprétables par le modèle dans les étapes suivantes.

    \item \textbf{Exécuter le modèle agent} : le modèle consiste en un réseau de neurones profond à deux têtes (architecture acteur-critique simplifiée) avec des connexions résiduelles qui prend les 1024 dimensions de l'Étape 1 et prédit :
    \begin{itemize}
        \item \textbf{Tête d'action (policy)} : un vecteur de taille 256 dimensions représentant l'embedding du prompt à sélectionner
        \item \textbf{Tête de valeur} : un scalaire estimant la récompense attendue pour guider l'apprentissage via l'estimation d'avantage
    \end{itemize}
    
    L'architecture améliorée est la suivante :
    \begin{itemize}
        \item \textbf{Projection d'entrée} : $1024 \rightarrow 512$ avec LayerNorm, GELU et Dropout (probabilité = 0.15, réduite de 0.2 pour améliorer le flux de gradient pendant les récompenses rares)
        \item \textbf{Blocs résiduels} : 4 blocs résiduels avec pré-normalisation (augmenté de 3 à 4 pour améliorer la capacité d'apprentissage de représentations). Chaque bloc implémente :
        \[
        \text{ResidualBlock}(x) = x + \text{MLP}(\text{LayerNorm}(x))
        \]
        où le MLP utilise un facteur d'expansion de 4 : $512 \rightarrow 2048 \rightarrow 512$ avec activation GELU. Cette architecture permet un meilleur flux de gradient et une capacité représentationnelle accrue.
        \item \textbf{Tête d'action} : $512 \rightarrow 512 \rightarrow 256$ avec normalisation L2 pour l'alignement avec l'index FAISS. Utilise un dropout réduit de 50\% dans la couche finale pour préserver la cohérence des prédictions.
        \item \textbf{Tête de valeur améliorée} : architecture séparée pour prévenir l'enchevêtrement avec la politique : $512 \rightarrow 256 \rightarrow 64 \rightarrow 1$ avec activation Sigmoid pour contraindre la sortie à [0, 1] correspondant à la plage de récompenses. Une projection intermédiaire supplémentaire améliore la capacité d'estimation.
    \end{itemize}
    
    Les stratégies d'exploration et de régularisation implémentées sont :
    \begin{itemize}
        \item \textbf{Exploration $\epsilon$-greedy consciente de la variété (\textit{manifold-aware})} : avec une probabilité $\epsilon$ fixe (paramètre configurable, par défaut $\epsilon = 0.0$ pour privilégier l'exploitation), le modèle peut sélectionner un vecteur d'exploration. Contrairement à l'approche naïve qui échantillonne aléatoirement sur l'hypersphère (produisant souvent des vecteurs hors de la variété des prompts valides), notre approche utilise des \textbf{centroïdes pré-calculés} par k-means sur les embeddings des prompts. L'exploration échantillonne parmi ces centroïdes avec un léger bruit local ($\sigma = 0.1$), garantissant que les vecteurs explorés restent sur la variété des prompts réels. Ces centroïdes sont sauvegardés dans \texttt{prompt\_centroids.pt} et chargés automatiquement au démarrage du modèle.
        
        \item \textbf{Bruit d'exploration} : ajout d'un bruit gaussien de $\sigma = 0.15$ (réduit de 0.2) à la prédiction durant l'entraînement pour améliorer la robustesse du modèle face aux variations et encourager l'exploration de l'espace des embeddings tout en évitant le surapprentissage. Après l'ajout du bruit, le vecteur est renormalisé en L2 pour maintenir la cohérence avec l'index FAISS qui stocke des vecteurs normalisés.
        
        \item \textbf{Normalisation adaptative des récompenses} : utilisation de statistiques courantes (moyenne et écart-type mobiles) calculées selon l'algorithme de Welford pour normaliser les récompenses et stabiliser l'apprentissage face à leur grande variance. L'écart-type est contraint à la plage [0.1, 0.5] pour éviter une normalisation trop agressive. La normalisation clipse les valeurs à [-3, 3] pour prévenir les valeurs extrêmes.
        
        \item \textbf{Pénalité de diversité avec décroissance} : suivi de la fréquence d'utilisation de chaque prompt avec un facteur de décroissance de 0.99 pour éviter une croissance infinie des pénalités. La pénalité utilise une transformation logarithmique douce pour un échelle plus modéré : 
        \[
        \text{penalty}_i = \alpha \cdot \log(1 + 10 \cdot \text{freq}_i), \quad \text{freq}_i = \frac{\text{count}_i}{\sum_j \text{count}_j}
        \]
        où $\alpha = 0.05$ (réduit de 0.1) est le coefficient de pénalité, et la pénalité maximale est plafonnée à 0.1.
    \end{itemize}

    \item \textbf{Mesurer la similarité} : entre la prédiction (vecteur 256 dimensions normalisé L2) et tous les vecteurs de la base de données vectorielle FAISS afin de sélectionner le prompt le plus proche de la prédiction par similarité cosinus et le retourner avec son audio et sa transcription associés. La base vectorielle contient les embeddings de 1000 prompts (phase de test) ou 63000 prompts (entraînement complet) réduits de 1024 à 256 dimensions par ACP.

    \item \textbf{Générer l'audio avec ZipVoice} : le système prend la transcription et le segment audio du prompt sélectionné à l'étape précédente ainsi que le texte cible et génère le nouveau segment audio par clonage de voix en \textit{one-shot}.

    \item \textbf{Mesurer le Character Error Rate (CER)} : en utilisant Whisper Large v3 \cite{whisper}, un modèle de reconnaissance automatique de la parole (\textit{Automatic Speech Recognition}, ASR). L'idée est d'obtenir la transcription du segment de référence provenant du corpus et du segment qui vient d'être généré par ZipVoice en utilisant la prédiction comme prompt. Le CER mesure la différence entre ces deux transcriptions caractère par caractère selon la distance d'édition de Levenshtein, avec une erreur plus basse indiquant une meilleure qualité de synthèse vocale.

    \item \textbf{Calculer la récompense pondérée} : comme le modèle s'entraîne selon une stratégie d'apprentissage par renforcement, il faut calculer et fournir une récompense pour guider l'agent vers de meilleurs choix. Si la récompense est élevée, l'agent va apprendre beaucoup d'informations de cet exemple et les poids du modèle vont se mettre à jour fortement dans la direction de cette action optimale. Si la récompense est faible, le modèle va à peine apprendre de cet exemple et son état va rester approximativement le même. 
    
    Dans notre cas, nous avons défini la récompense par la formule : 
    \[
    R = \max(0, 1 - \mathrm{CER})
    \]
    où $R$ désigne la récompense (\textit{Reward}), et $\mathrm{CER}$ correspond au taux d'erreur de caractères. Cette formule s'interprète comme suit :
    \begin{itemize}
      \item Si $\mathrm{CER} = 0.0$ (transcription parfaite) $\rightarrow$ $R = 1.0$
      \item Si $\mathrm{CER} = 0.5$ (50\% d'erreurs) $\rightarrow$ $R = 0.5$
      \item Si $\mathrm{CER} \geq 1.0$ (erreurs majeures) $\rightarrow$ $R = 0.0$
    \end{itemize}
    Il s'agit d'une transformation linéaire inverse qui convertit le CER (où plus bas = meilleur) en récompense (où plus haut = meilleur).

    \item \textbf{Mise à jour du modèle avec Policy Gradient et apprentissage contrastif} : dans cette étape finale, le modèle effectue une mise à jour combinant un gradient de politique de type REINFORCE et l'apprentissage contrastif InfoNCE. Cette approche représente une amélioration majeure par rapport à la \textit{Reward-Weighted Regression} (RWR) simple, permettant un assignement de crédit approprié et un apprentissage à partir d'exemples positifs et négatifs.
    
    \textbf{Estimation d'avantage avec baseline} : le modèle calcule l'avantage comme la différence entre la récompense observée et une baseline mobile :
    \[
    A = R - R_{\text{baseline}}
    \]
    où $R_{\text{baseline}}$ est mise à jour par moyenne mobile exponentielle : $R_{\text{baseline}} \leftarrow 0.99 \cdot R_{\text{baseline}} + 0.01 \cdot R$. Les avantages sont normalisés par l'écart-type du batch pour stabiliser l'apprentissage.
    
    La fonction de perte combinée est :
    \[
    \mathcal{L}_{\text{total}} = \mathcal{L}_{\text{PG}} + 0.3 \cdot \mathcal{L}_{\text{InfoNCE}} + 0.5 \cdot \mathcal{L}_{\text{value}} + \mathcal{L}_{\text{value-reg}} - \alpha \cdot H(\text{prédiction}) + \beta \cdot \text{Penalty}_{\text{div}}
    \]
    
    Les composantes de cette fonction sont :
    \begin{itemize}
        \item \textbf{Perte de gradient de politique (Policy Gradient)} : 
        \[
        \mathcal{L}_{\text{PG}} = -\mathbb{E}\left[\log(\text{sim}(\text{prédiction}, \text{action})) \cdot A_{\text{norm}}\right]
        \]
        où $\text{sim}$ est la similarité cosinus entre la prédiction du modèle et l'action sélectionnée, et $A_{\text{norm}}$ est l'avantage normalisé. Cette approche de type REINFORCE permet au modèle d'apprendre à maximiser la similarité avec les actions à haute récompense et à la minimiser pour les actions à faible récompense.
        
        \item \textbf{Perte contrastive InfoNCE} : lorsque le batch contient suffisamment d'échantillons, les actions sont séparées en positives (récompense $\geq$ médiane) et négatives (récompense $<$ médiane). La perte InfoNCE attire les prédictions vers les actions positives et les repousse des actions négatives :
        \[
        \mathcal{L}_{\text{InfoNCE}} = -\log\frac{\exp(\text{sim}(\text{pred}, \text{pos}^+)/\tau)}{\exp(\text{sim}(\text{pred}, \text{pos}^+)/\tau) + \sum_j \exp(\text{sim}(\text{pred}, \text{neg}_j)/\tau)}
        \]
        où $\tau = 0.1$ est la température. Cette approche permet au modèle d'apprendre non seulement ce qu'il faut imiter, mais aussi ce qu'il faut éviter.
        
        \item \textbf{Perte de valeur avec régularisation} : 
        \[
        \mathcal{L}_{\text{value}} = \text{MSE}(\text{valeur\_prédite}, R) - 0.01 \cdot \text{Var}(\text{valeur\_prédite})
        \]
        La perte MSE entraîne la tête de valeur à estimer précisément les récompenses futures. Le terme de variance négative prévient l'effondrement de la tête de valeur (où toutes les prédictions deviennent identiques), permettant un calcul d'avantage robuste pour guider l'exploration.
        
        \item \textbf{Régularisation par entropie} : encourage la diversité des prédictions avec coefficient $\alpha = 0.1$ (valeur par défaut, ajustable). L'entropie est approximée par :
        \[
        H = \text{std}(\text{prédiction}) + 0.5 \cdot \text{mean}(\text{pdist}(\text{prédictions}))
        \]
        où $\text{std}$ mesure la variation au sein de chaque prédiction et $\text{pdist}$ mesure la diversité entre prédictions du batch, évitant l'effondrement de mode (\textit{mode collapse}).
        
        \item \textbf{Pénalité de diversité} : $\text{Penalty}_{\text{div}}$ avec poids $\beta = 0.05$ (réduit de 0.02) pénalise la surexploitation de certains prompts basée sur leur fréquence d'utilisation historique avec décroissance temporelle.
    \end{itemize}
    
    \textbf{Stratégie d'entraînement avec Replay Buffer} : le modèle utilise un \textit{replay buffer} de capacité maximale 5000 expériences (augmentée de 1000), implémenté avec une structure \textit{deque} qui élimine automatiquement les expériences les plus anciennes. À chaque itération, le modèle échantillonne aléatoirement des mini-batches de taille 64 depuis ce buffer pour l'entraînement sur 10 époques (réduit de 20 pour un entraînement plus rapide tout en préservant la qualité). 
    
    L'échantillonnage utilise une stratégie consciente de la diversité (\textit{diversity-aware sampling}) : parmi un ensemble de candidats (4× la taille du batch), le système sélectionne un sous-ensemble diversifié en maximisant itérativement la distance minimale aux échantillons déjà sélectionnés. Cette approche permet de :
    \begin{itemize}
        \item Réutiliser efficacement les expériences passées
        \item Briser les corrélations temporelles entre exemples consécutifs
        \item Prévenir l'oubli catastrophique
        \item Améliorer la stabilité et l'efficacité de l'échantillonnage
        \item Assurer une couverture uniforme de l'espace des actions
    \end{itemize}
    
    Les gradients sont clippés avec une norme maximale de 1.0 pour assurer la stabilité de l'apprentissage, et l'optimiseur utilisé est Adam avec un taux d'apprentissage de $1 \times 10^{-4}$ (réduit par rapport à la version initiale de $5 \times 10^{-4}$ pour améliorer la stabilité), avec option d'utiliser un scheduler d'annealing cosinus pour réduire progressivement le taux d'apprentissage selon :
    \[
    \eta_t = \eta_{\min} + \frac{1}{2}(\eta_{\max} - \eta_{\min})\left(1 + \cos\left(\frac{t}{T_{\max}}\pi\right)\right)
    \]
    où $\eta_{\max} = 1 \times 10^{-4}$, $\eta_{\min} = 1 \times 10^{-5}$, et $T_{\max}$ est le nombre total de pas d'entraînement.
    
    \textbf{Sauvegarde et évaluation} : le système sauvegarde automatiquement des checkpoints tous les 100 phrases traitées, incluant à la fois le modèle et le replay buffer. Tous les 50 phrases, une évaluation est effectuée sur un ensemble de validation distinct (10 premières phrases de \texttt{sentences\_val.txt}) sans exploration ($\epsilon = 0$) pour mesurer la performance réelle du modèle. Les métriques d'entraînement (numéro de phrase, CER, récompense, similarité cosinus) sont enregistrées dans \texttt{training\_progress.csv} pour analyse et visualisation de la convergence.
    
\end{enumerate}
